% Report: calibration of jump and stochastic-volatility models to Bitcoin options, in the browser.
% Build: in docs/report, pdflatex report.tex (twice), then move report.pdf to ../../REPORT.pdf.
% The figures read data/*.dat, extracted from the application itself (day 2026-08-26) and from jeux.js.
\documentclass[11pt,a4paper]{article}
\usepackage[a4paper,margin=2.4cm]{geometry}
\usepackage[T1]{fontenc}
\usepackage[utf8]{inputenc}
\usepackage{lmodern}
\usepackage{microtype}
\usepackage{amsmath,amssymb,amsthm,mathtools,bm}
\usepackage{booktabs}
\usepackage{enumitem}
\usepackage{xcolor}
\usepackage{pgfplots}
\pgfplotsset{compat=1.17}
\usepgfplotslibrary{dateplot,groupplots}
\usepackage{caption}
\usepackage[colorlinks=true,linkcolor=blue!50!black,citecolor=blue!50!black,urlcolor=blue!50!black]{hyperref}
\usepackage{cleveref}
\hypersetup{pdftitle={Calibrating jump and stochastic-volatility models to Bitcoin options},pdfauthor={Mohammed-Yassine Barnicha}}

\newtheorem{proposition}{Proposition}
\theoremstyle{definition}
\newtheorem{definition}{Definition}
\theoremstyle{remark}
\newtheorem{remark}{Remark}

\newcommand{\E}{\mathbb{E}}
\newcommand{\Q}{\mathbb{Q}}
\newcommand{\N}{\mathcal{N}}
\newcommand{\dd}{\,\mathrm{d}}
\newcommand{\ii}{\mathrm{i}}
\pgfplotsset{every axis/.append style={grid=major,grid style={gray!25},scaled ticks=false,
  tick label style={font=\footnotesize},label style={font=\footnotesize},title style={font=\footnotesize},
  legend style={font=\scriptsize,draw=none,fill=white,fill opacity=0.8,text opacity=1}}}

\title{Calibrating Jump and Stochastic-Volatility Models\\to Bitcoin Options, in the Browser\\[4pt]
  \large Model composition, Fourier pricing, loss design, optimisation and validation}
\author{Mohammed-Yassine Barnicha}
\date{October 2026}

\begin{document}
\maketitle

\begin{abstract}
\noindent This report documents an application that prices, calibrates and analyses Deribit Bitcoin options with 28
affine models, entirely in a web browser. Each model is the composition of a variance block (Black--Scholes, Heston, double
Heston, Heston with variance jumps) and a jump law (none, Merton, bi-lognormal, Kou, CGMY, Variance Gamma, NIG); all are
priced through their characteristic function by the COS and PROJ methods, re-implemented in JavaScript and checked against
a Python reference to $6\cdot10^{-6}$\,\$. We derive the characteristic functions and their martingale constraints, the two
Fourier methods with their error control, the loss functions and the vega weighting used in production, the optimisers (classic
and adaptive Nelder--Mead, differential evolution), the multi-day calibration, the model Greeks, the gamma exposure and the
Monte Carlo engine for exotics. On 26 August 2026 the Bates model fits 222 out-of-the-money options over six maturities with
an implied-volatility RMSE of $37.3$\,bp. The application is at \url{https://yassinebarnicha.github.io/calibration-btc-svjd/}.
\end{abstract}

\tableofcontents

%======================================================================
\section{Introduction}
Bitcoin options trade mostly on Deribit, with expiries from one day to more than a year and a smile that differs from equity
markets: the skew changes sign over time, the wings are steep at short maturities, and jumps of several percent are routine.
A model calibrated to this surface must combine a stochastic variance (term structure, smile dynamics) and jumps (short-dated
wings); which combination works best is an empirical question, and answering it requires calibrating many models on many days.

The application makes that experiment interactive and portable. Its design goals are:
\begin{itemize}[nosep]
  \item \emph{no installation}: a static page, the pricers in JavaScript, every day of data in a file loaded on demand;
  \item \emph{transparency}: each model shows its dynamics, its constraints and its loss; every numerical choice is adjustable;
  \item \emph{speed}: a full evaluation of the loss on a production day takes about $1.3$\,ms, so sliders reprice the whole chain
        live and an optimiser runs about $650$ evaluations per second;
  \item \emph{reproducibility}: the default settings reproduce the production calibration book exactly.
\end{itemize}

%======================================================================
\section{Market data}\label{sec:data}
\subsection{Instruments and conventions}
Deribit options are European, written on the BTC/USD index $X$ and quoted in BTC (inverse options): the premium of a call is
$C^{\mathrm{BTC}}=C^{\mathrm{USD}}/X$. For each expiry $T$ a future (or, for daily expiries, an implied forward) gives the
forward $F_T$. Prices are converted to dollars, $P=P^{\mathrm{BTC}}X$, and the implied rate of the maturity is
$r=\ln(F_T/X)/T$. Options are then priced with Black-76 on the forward, $P=e^{-rT}\,\mathrm{B76}(F,K,T,\sigma)$, and implied
volatilities are obtained by inverting Black-76 (safeguarded Newton, bracket $[10^{-6},5]$, tolerance $10^{-10}$ on the price).
Time to expiry is counted to the hour up to the 08:00\,UTC expiry, which matters for the shortest maturities.

\subsection{The production book}
The target price of an option is its mark if $\mathrm{bid}\le\mathrm{mark}\le\mathrm{ask}$, otherwise the mid. The production
filters keep out-of-the-money options (puts with $K<F$, calls with $K>F$), quoted on both sides, with a relative spread
$(\mathrm{ask}-\mathrm{bid})/\mathrm{mid}<20\,\%$ and a target price above $0.005$\,BTC, on the last six maturities beyond ten days.
On 26 August 2026 this leaves 222 of the 914 options of the chain. Every filter can be changed in the application, including
maturity-dependent thresholds (a $0.004$\,BTC option is a normal wing price at one week, an illiquid tail at one year).

\subsection{Data pipeline}
A day is a JavaScript file \texttt{jours/<date>.js} holding the whole chain in columns (strikes, maturities, $F$, $r$, bid,
ask, mark, their implied volatilities, open interest) and the parameter sets calibrated that day. The page loads it with a
\texttt{<script>} element, so it also works from the local disk. A GitHub Action downloads the chain from the public Deribit API
every morning after the 08:00\,UTC expiry, adds the day and rebuilds the page; the ``Today'' button does the same from the
browser. The data cover 270 days since January 2024 and 529 calibrated parameter sets.

%======================================================================
\section{The models}\label{sec:models}
\subsection{Composition}
Under the risk-neutral measure the future is a martingale. All models write
\begin{equation}\label{eq:compo}
  \ln\frac{F_T}{F_0}=X^{c}_T+X^{J}_T-\omega T,
\end{equation}
where $X^c$ is the continuous part driven by the variance block, $X^J$ a pure-jump Lévy process independent of it, and $\omega$
the compensator that makes $F$ a martingale.
\begin{proposition}\label{prop:compo}
Let $\psi_J$ be the Lévy exponent of $X^J$, $\E[e^{\ii uX^J_T}]=e^{T\psi_J(u)}$, and $A_{\mathrm{var}}(u,T)$ the log characteristic
function of $X^c_T$, with $\E[e^{X^c_T}]=1$. Then
\begin{equation}
  \varphi_T(u)=\E\big[e^{\ii u\ln(F_T/F_0)}\big]=\exp\Big(A_{\mathrm{var}}(u,T)+T\psi_J(u)-\ii u\,\omega T\Big),\qquad \omega=\psi_J(-\ii),
\end{equation}
and $F$ is a martingale iff $\psi_J(-\ii)$ is finite.
\end{proposition}
\begin{proof}
Independence factorises the characteristic function. $\E[F_T/F_0]=\E[e^{X^c_T}]\,\E[e^{X^J_T}]\,e^{-\omega T}
=e^{T\psi_J(-\ii)-\omega T}$, which equals one for all $T$ iff $\omega=\psi_J(-\ii)$; the Markov property of $(F,v)$ then gives
the martingale property at all intermediate dates.
\end{proof}
The page assembles the pricer object from the two blocks: a block that is absent leaves no parameter at all (``no jumps''
removes $\lambda$ from the sliders, the bounds and the exports, it is not $\lambda=0$).

\subsection{Variance blocks}
\paragraph{Black--Scholes.} $A_{\mathrm{var}}=-\tfrac12\sigma^2T(u^2+\ii u)$.
\paragraph{Heston.} $\dd F_t/F_t=\sqrt{v_t}\dd W_t$, $\dd v_t=\kappa(\theta-v_t)\dd t+\sigma\sqrt{v_t}\dd Z_t$, $\dd\langle W,Z\rangle_t=\rho\dd t$.
The exponent is affine in $v_0$, $A_{\mathrm{var}}=C(u,T)+D(u,T)v_0$, with the Riccati solution written in the form that avoids
the branch-cut discontinuity of the complex logarithm (Albrecher et al.\ 2007):
\begin{gather}
  d=\sqrt{(\rho\sigma\ii u-\kappa)^2+\sigma^2(u^2+\ii u)},\qquad g=\frac{\kappa-\rho\sigma\ii u-d}{\kappa-\rho\sigma\ii u+d},\\
  C=\frac{\kappa\theta}{\sigma^2}\Big[(\kappa-\rho\sigma\ii u-d)T-2\ln\frac{1-ge^{-dT}}{1-g}\Big],\qquad
  D=\frac{\kappa-\rho\sigma\ii u-d}{\sigma^2}\,\frac{1-e^{-dT}}{1-ge^{-dT}} .
\end{gather}
The Feller condition $2\kappa\theta\ge\sigma^2$ keeps $v$ strictly positive; the optimiser can enforce it.
\paragraph{Double Heston.} Two independent CIR factors (Christoffersen, Heston and Jacobs 2009): the exponents add. A fast factor
(large $\kappa_1$) shapes the short end, a slow one the long end; Feller is checked factor by factor.
\paragraph{Heston with variance jumps.} $\dd v_t=\kappa(\theta-v_t)\dd t+\sigma\sqrt{v_t}\dd Z_t+\dd J^v_t$ with $J^v$ compound
Poisson of intensity $\lambda_v$ and exponential sizes of mean $\mu_v$ (Duffie, Pan and Singleton 2000). The exponent gains
\begin{equation}
  \lambda_v\int_0^T\Big(\frac{1}{1-\mu_vD(u,t)}-1\Big)\dd t,
\end{equation}
computed by Simpson's rule on 64 intervals. The exponential law is what keeps the problem affine with a one-dimensional
integral; a lognormal size would require a numerical integral inside the Fourier inversion, too slow for live use.

\subsection{Jump laws}
\Cref{tab:jumps} lists the Lévy exponents. Their constraints are those of the existence of $\psi_J(-\ii)$, i.e.\ of $\E[e^{X^J_1}]$.
\begin{table}[ht]\centering\small
\renewcommand{\arraystretch}{1.35}
\begin{tabular}{@{}lll@{}}\toprule
law & $\psi_J(u)$ & constraint\\\midrule
Merton & $\lambda\big(e^{\ii u\mu_J-\sigma_J^2u^2/2}-1\big)$ & $\lambda\ge0$, $\sigma_J>0$\\
bi-lognormal & sum of two Merton exponents & $\lambda_1,\lambda_2\ge0$ \\
Kou & $\lambda\Big(\dfrac{p\eta_1}{\eta_1-\ii u}+\dfrac{(1-p)\eta_2}{\eta_2+\ii u}-1\Big)$ & $\eta_1>1$, $\eta_2>0$\\
CGMY & $C\,\Gamma(-Y)\big[(M-\ii u)^Y-M^Y+(G+\ii u)^Y-G^Y\big]$ & $M>1$, $G>0$, $Y<2$, $Y\notin\{0,1\}$\\
Variance Gamma & $-\dfrac1\nu\ln\Big(1-\ii u\theta\nu+\tfrac12\sigma^2\nu u^2\Big)$ & $1-\theta\nu-\tfrac12\sigma^2\nu>0$\\
NIG & $\delta\Big(\sqrt{\alpha^2-\beta^2}-\sqrt{\alpha^2-(\beta+\ii u)^2}\Big)$ & $\alpha>|\beta+1|$, $\delta>0$\\\bottomrule
\end{tabular}
\caption{Lévy exponents of the jump blocks (Merton 1976; Kou 2002; Carr, Geman, Madan and Yor 2002; Madan, Carr and Chang
1998; Barndorff-Nielsen 1997).}\label{tab:jumps}
\end{table}

\begin{remark}[Identification]
Several compositions have flat directions that no optimiser can resolve from one day of data: the two components of the
bi-lognormal law are interchangeable unless their intensities live in disjoint ranges (production bounds $\lambda_1\in[1,300]$,
$\lambda_2\in[0,1]$); in Heston, $v_0$ and $\theta$ trade off at short horizons and $\kappa$ and $\sigma$ along a ridge; jump
intensity and jump size trade off for the short-dated wings. The page makes these visible: parameters can be locked, swapped
(``Swap components 1 and 2'') or freed per day when several days are calibrated together (\cref{sec:stack}).
\end{remark}

%======================================================================
\section{Fourier pricing}\label{sec:fourier}
Both methods need only $\varphi_T$; the first two cumulants $c_1,c_2$ of $\ln(F_T/F_0)$ are obtained by central differences of
$\ln\varphi_T$ at $u=0$ (step $10^{-5}$), and the fourth $c_4$ in closed form for each jump family.

\subsection{The COS method}
Fang and Oosterlee (2008) expand the density of $x=\ln(F_T/F_0)$ on a truncated interval $[a,b]$ in a cosine series whose
coefficients are read directly from the characteristic function:
\begin{equation}
  f(x)\approx\frac{2}{b-a}{\sum_{k=0}^{N-1}}{}'\;\mathrm{Re}\Big[\varphi_T\big(u_k\big)e^{-\ii u_ka}\Big]\cos\big(u_k(x-a)\big),\qquad u_k=\frac{k\pi}{b-a},
\end{equation}
where $\sum'$ halves the first term. Integrating a payoff against this expansion uses the closed-form cosine moments
$\chi_k,\psi_k$ of $e^x$ and $1$ on sub-intervals. Puts are priced (their payoff is bounded) and calls follow from put--call
parity $C-P=e^{-rT}(F-K)$, which removes the cancellation error of deep in-the-money calls. The implementation uses
$[a,b]=c_1\pm L\sqrt{c_2}$ with $L=24$ and $N=512$ by default. Since all options of a maturity share $\varphi_T(u_k)$, these
$N$ values are computed once per maturity; the series is cut where a bound on its remainder falls below $10^{-10}$\,\$:
long maturities, whose characteristic function decays fast, need only a fraction of the terms.

\paragraph{Errors.} The error has two parts: truncation of the density to $[a,b]$, controlled by $L$ and the tails, and the
series truncation, which decays exponentially for densities that are smooth on $[a,b]$. COS needs a valid second cumulant and
loses accuracy when the jump part dominates the variance at very short maturities, where the density is spiky.

\subsection{The PROJ method}
Kirkby (2015) projects the density on a basis of linear B-splines on a uniform grid of $N$ points and spacing $\Delta$. The
projection coefficients are obtained by one FFT of $\varphi_T$ multiplied by the Fourier transform of the basis, and the payoff
is integrated against the basis in closed form (with a three-point Gauss rule for the boundary cell). The grid half-width is
\begin{equation}
  \alpha=\max\Big(L\sqrt{|c_2|+\sqrt{|c_4|}},\ 1.15\max_K|\ln(K/F)|+c_1\Big),
\end{equation}
$L=20$, $N=2^{11}$ by default. The error is $O(\Delta^2)$ for smooth densities and the method stays robust when $|\mu_J|$,
$\sigma_J$ or $\lambda$ are large, which is why it serves as the reference for exotic controls and the volatility surface. COS
is about three times faster for the calibration loop.

\subsection{Validation of the pricers}
\begin{itemize}[nosep]
  \item Put--call parity at the money: $C(K=F)-P(K=F)=0.0000$\,bp on the 28 compositions (martingale property).
  \item Black--Scholes recovered to within $0.005$\,bp; COS and PROJ equal on a full smile.
  \item JavaScript pricers equal to the Python reference pricers to $6\cdot10^{-6}$\,\$.
  \item Successive optimisations of the code checked against the previous version on all compositions, COS and PROJ: equal to
        rounding noise ($4.5\cdot10^{-8}$ relative).
\end{itemize}

%======================================================================
\section{The calibration problem}\label{sec:calib}
\subsection{Loss functions}
For options $i$ with target price $P_i$, implied volatility $\sigma_i$ and Black-76 vega $\mathcal V_i$ (at the target volatility),
and model values $\hat P_i(\Theta)$, $\hat\sigma_i(\Theta)$, the loss is a weighted mean of squared errors
\begin{equation}
  \mathcal L(\Theta)=\frac{\sum_iw_i\,e_i(\Theta)^2}{\sum_iw_i},
\end{equation}
with five error measures (price, relative price, price over the bid--ask range, implied volatility in bp, and a combination of
price and volatility with a weight $k$) and five weightings (none, raw $1/\mathcal V_i^2$, $1/\mathcal V_i^2$ normalised per
maturity to mean one or to sum one, equal weight per maturity). The page displays the formula being minimised.
\begin{proposition}[Vega weighting]
To first order, $P_i-\hat P_i\approx\mathcal V_i\,(\sigma_i-\hat\sigma_i)$, so $\sum_i(P_i-\hat P_i)^2/\mathcal V_i^2\approx\sum_i(\sigma_i-\hat\sigma_i)^2$.
\end{proposition}
The price loss with $1/\mathcal V^2$ weights therefore approximates the implied-volatility loss without inverting Black-76 at each
evaluation. The production loss, \texttt{inv\_vega2\_equal}, normalises the weights to sum one per maturity so that every
maturity counts equally whatever its number of options; with raw weights, short maturities (small vega) dominate. A floor of
$10^{-4}$ on the vega avoids infinite weights. The IV and price RMSE are reported whatever the loss.

\subsection{Search space}
Each free parameter is mapped to a slider $t\in[0,1]$ inside its bounds $[\ell,h]$: $\Theta=\ell(h/\ell)^t$ for positive scale
parameters (log scale), $\Theta=\ell+t(h-\ell)$ for $\rho$ and $\mu_J$, and $\Theta=\ell+0.01\,(e^{t\ln(1+(h-\ell)/0.01)}-1)$ for
intensities, which can reach zero exactly. The optimisers work in $t$; leaving $[0,1]$, or violating Feller when it is enforced,
gives $\mathcal L=+\infty$. The default bounds are very wide; the production bounds of each model can be loaded.

\subsection{Optimisers}
\begin{itemize}
  \item \emph{Classic Nelder--Mead} (reflection, expansion, contraction, shrink coefficients $1,2,\tfrac12,\tfrac12$), from the
        current parameters, initial step $0.04$ in slider units, stopping when $(f_{\max}-f_{\min})/f_{\min}\le10^{-6}$ or after
        2500 evaluations.
  \item \emph{Adaptive Nelder--Mead} (Gao and Han 2012): coefficients $1$, $1+2/n$, $\tfrac34-\tfrac1{2n}$, $1-\tfrac1n$ in dimension
        $n$, which keeps the simplex from degenerating when many parameters are free.
  \item \emph{Differential evolution} (Storn and Price 1997): a population of $\text{popsize}\times n$ points initialised by Latin
        hypercube, mutation $v=x_{r_1}+F(x_{r_2}-x_{r_3})$ (rand1bin) or $v=x_{\text{best}}+F(x_{r_1}-x_{r_2})$ (best1bin) with $F$
        drawn in $[F_{\min},F_{\max}]$ (dithering), binomial crossover with rate $CR$ and greedy selection; it stops when the
        standard deviation of the population's losses falls below $\text{atol}+\text{tol}\,|\text{mean}|$. The production settings
        are rand1bin, $F\in[0.7,1.5]$, $CR=0.6$, popsize 25, tol $10^{-7}$; the current parameters can be injected in the population.
  \item \emph{DE then Nelder--Mead}: the Python production chain (global then local), with Nelder--Mead in place of SLSQP.
\end{itemize}
Same seed, same settings, same data: same result. The optimisers run in 12\,ms slices between two frames, so the page stays
fluid (60 frames per second) during a calibration.

\subsection{Several days at once}\label{sec:stack}
A contract observed on two consecutive days gives two maturities $T$ and $T-1/365$: stacking $N$ days multiplies the points of
the term structure by $N$ without new data. Parameters are shared by all days, which is the identifying assumption; the state
variables ($v_0$, $\theta$, $\rho$, jump intensities) can be freed day by day. The weighting can give each maturity, or each day,
a total weight of one.

%======================================================================
\section{Analysis tools}\label{sec:tools}
\subsection{Model Greeks}
Greeks are computed by central finite differences on the model's own Fourier pricer, with the parameters fixed: $\Delta$ and
$\Gamma$ by moving $F$ by $\pm0.1\,\%$, vega by scaling the instantaneous volatility $\sigma_0$ by $\pm0.1\,\%$ ($\sigma$ in
Black--Scholes, $\sqrt{v_0}$ in Heston, $\sqrt{v_{0,1}+v_{0,2}}$ in double Heston), theta by moving the maturity by one day on each
side. With the Black--Scholes block, these differences recover the Black-76 closed forms to the decimals shown for $\Delta$,
$\Gamma$ and vega, and within $0.1\,\%$ for $\Theta$. The gap between the model delta and the Black-76 delta at the market
volatility is the smile effect. With its parameters fixed, the model's implied volatility is a function of moneyness, so when
$F$ moves the option slides along the smile: $\Delta_{\text{model}}=\Delta_{\text{B76}}+\mathcal V\,\partial\sigma_{\text{imp}}/\partial F$.
The sign of the correction follows the slope of the smile, and it is the extra hedge the model recommends.

\subsection{Volatility surface}
The model surface is priced by PROJ on a strike grid per maturity (25 to 61 strikes) and inverted with Black-76. It is shown
either on the calibrated range only or over all the day's options, where it is an extrapolation of what the model has seen.

\subsection{Gamma exposure}
The gamma exposure of option $i$ is $\mathrm{GEX}_i=s_i\cdot\mathrm{OI}_i\cdot\Gamma_i\cdot F^2\cdot1\,\%$, the dollar delta to trade
to hedge a $+1\,\%$ move of the future, with open interest in BTC. The ``market makers'' convention takes $s_i=+1$ for calls and
$-1$ for puts (dealers assumed long the public's calls and short its puts). The profile recomputes the total GEX for a future
moved from $-20\,\%$ to $+20\,\%$ with the volatilities of each strike frozen; its zero is the \emph{flip level}. It is a convention:
Deribit does not publish who holds what.

\subsection{Exotic options by Monte Carlo}
Paths of the future are simulated with the active model. The CIR variance uses the quadratic--exponential scheme of Andersen
(2008): given $v_t$, the conditional mean $m$ and variance $s^2$ of $v_{t+\Delta}$ are matched by $a(b+Z)^2$ when
$s^2/m^2\le1.5$ and by a mixture of a mass at zero and an exponential law otherwise. Lognormal, Kou, Variance Gamma (gamma
subordinator) and NIG (inverse Gaussian subordinator) jumps are simulated exactly per step and compensated, so the future stays
a martingale; CGMY has no simple exact simulation and is excluded. Averages, extrema and barriers are monitored at each step.
American options use Longstaff--Schwartz with the basis $1,x,x^2,x^3$, $x=F/K$, on in-the-money paths. The European with the same
strike is simulated on the same paths and compared with its Fourier price: the gap measures the discretisation bias, and the
European serves as a control variate, $\hat Y_{\mathrm{cv}}=\hat Y-b^\star(\hat C-C_{\mathrm{Fourier}})$ with
$b^\star=\operatorname{Cov}(Y,C)/\operatorname{Var}(C)$. On six model families the simulated European falls within two standard
errors of the Fourier price; the early-exercise premium is zero for $r=0$ and positive for $r>0$, as it must be.

%======================================================================
\section{Results}\label{sec:results}
\subsection{A calibrated day}
\Cref{tab:bates,fig:smiles} show the Bates model (Heston + Merton) on 26 August 2026, with its production calibration
(differential evolution then local search, COS, price loss with \texttt{inv\_vega2\_equal} weights).

\begin{table}[ht]\centering\small
\begin{tabular}{@{}lcccccccc@{}}\toprule
 & $\kappa$ & $\theta$ & $\sigma$ & $\rho$ & $v_0$ & $\lambda$ & $\mu_J$ & $\sigma_J$\\\midrule
value & $46.03$ & $0.1334$ & $2.491$ & $0.139$ & $0.1630$ & $0.149$ & $-0.460$ & $0.783$\\\bottomrule
\end{tabular}
\caption{Bates parameters on 2026-08-26: IV RMSE $37.3$\,bp and price RMSE $44.0$\,\$ on 222 options, Feller margin
$2\kappa\theta-\sigma^2=6.08$.}\label{tab:bates}
\end{table}

The fit has a clear reading. The variance reverts very fast ($\kappa=46$, a half-life of $\ln2/\kappa\approx5.5$ days) towards
$\sqrt\theta=36.5\,\%$ from $\sqrt{v_0}=40.4\,\%$, with a large volatility of variance: the term structure of the smile is short
and steep. The spot--volatility correlation is \emph{positive} ($\rho=0.14$), the opposite of equity indices: on that day upward
moves were expected to come with higher volatility, which lifts the call wing. The jump component is a rare crash
($\lambda=0.15$ per year, mean size $e^{-0.46}-1\approx-37\,\%$, dispersion $0.78$), which carries the put wing at long maturities.

\begin{figure}[ht]\centering
\begin{tikzpicture}
\begin{groupplot}[group style={group size=3 by 2,horizontal sep=1.3cm,vertical sep=2.1cm},width=0.36\textwidth,height=4.6cm,
  xlabel={$\ln(K/F)$},xticklabel style={/pgf/number format/fixed}]
\nextgroupplot[title={2026-09-11 (16 d, 38 bp)},ylabel={IV (\%)},legend to name=leg,legend columns=3]
  \addplot[only marks,mark=*,mark size=1pt,black] table[x=k,y=mkt]{data/smile1.dat}; \addlegendentry{market (target)}
  \addplot[red!80!black,thick] table[x=k,y=mod]{data/smile1.dat}; \addlegendentry{Bates}
  \addplot[only marks,mark=-,gray] table[x=k,y=bid]{data/smile1.dat}; \addlegendentry{bid / ask IV}
  \addplot[only marks,mark=-,gray,forget plot] table[x=k,y=ask]{data/smile1.dat};
\nextgroupplot[title={2026-09-25 (30 d, 24 bp)}]
  \addplot[only marks,mark=*,mark size=1pt,black] table[x=k,y=mkt]{data/smile2.dat};
  \addplot[red!80!black,thick] table[x=k,y=mod]{data/smile2.dat};
  \addplot[only marks,mark=-,gray] table[x=k,y=bid]{data/smile2.dat}; \addplot[only marks,mark=-,gray] table[x=k,y=ask]{data/smile2.dat};
\nextgroupplot[title={2026-10-30 (65 d, 45 bp)}]
  \addplot[only marks,mark=*,mark size=1pt,black] table[x=k,y=mkt]{data/smile3.dat};
  \addplot[red!80!black,thick] table[x=k,y=mod]{data/smile3.dat};
  \addplot[only marks,mark=-,gray] table[x=k,y=bid]{data/smile3.dat}; \addplot[only marks,mark=-,gray] table[x=k,y=ask]{data/smile3.dat};
\nextgroupplot[title={2026-12-25 (121 d, 45 bp)},ylabel={IV (\%)}]
  \addplot[only marks,mark=*,mark size=1pt,black] table[x=k,y=mkt]{data/smile4.dat};
  \addplot[red!80!black,thick] table[x=k,y=mod]{data/smile4.dat};
  \addplot[only marks,mark=-,gray] table[x=k,y=bid]{data/smile4.dat}; \addplot[only marks,mark=-,gray] table[x=k,y=ask]{data/smile4.dat};
\nextgroupplot[title={2027-03-26 (212 d, 32 bp)}]
  \addplot[only marks,mark=*,mark size=1pt,black] table[x=k,y=mkt]{data/smile5.dat};
  \addplot[red!80!black,thick] table[x=k,y=mod]{data/smile5.dat};
  \addplot[only marks,mark=-,gray] table[x=k,y=bid]{data/smile5.dat}; \addplot[only marks,mark=-,gray] table[x=k,y=ask]{data/smile5.dat};
\nextgroupplot[title={2027-06-25 (303 d, 32 bp)}]
  \addplot[only marks,mark=*,mark size=1pt,black] table[x=k,y=mkt]{data/smile6.dat};
  \addplot[red!80!black,thick] table[x=k,y=mod]{data/smile6.dat};
  \addplot[only marks,mark=-,gray] table[x=k,y=bid]{data/smile6.dat}; \addplot[only marks,mark=-,gray] table[x=k,y=ask]{data/smile6.dat};
\end{groupplot}
\end{tikzpicture}

\pgfplotslegendfromname{leg}
\caption{The six production maturities of 2026-08-26 and the calibrated Bates smiles; the IV RMSE of each maturity is in its
title. Data exported from the application.}\label{fig:smiles}
\end{figure}

\subsection{Calibration history}
\Cref{fig:history} shows the best IV RMSE of each day among the calibrated sets embedded in the application, for Bates and for
the bi-lognormal Bates model. The median is $35.9$\,bp for Bates (61 days) and $37.7$\,bp for the bi-lognormal model (73 days):
the second jump component rarely pays for its three extra parameters on these days, but it removes the occasional failures of
a single jump size (the worst Bates day is above $1000$\,bp, a failed local search, against $114$\,bp for the bi-lognormal model).
The median correlation is slightly positive for both models.

\begin{figure}[ht]\centering
\begin{tikzpicture}
\begin{axis}[date coordinates in=x,width=0.92\textwidth,height=6cm,ymin=0,ymax=150,ylabel={best IV RMSE of the day (bp)},
  xticklabel={\year-\month},xtick={2024-01-01,2024-07-01,2025-01-01,2025-07-01,2026-01-01,2026-07-01},
  xticklabel style={font=\scriptsize},legend pos=north east,unbounded coords=jump,restrict y to domain=0:150]
  \addplot+[only marks,mark=*,mark size=1.2pt] table[x=date,y=rmse]{data/histoire_heston-logn.dat}; \addlegendentry{Bates (Heston + Merton)}
  \addplot+[only marks,mark=square*,mark size=1.2pt] table[x=date,y=rmse]{data/histoire_heston-biln.dat}; \addlegendentry{Heston + bi-lognormal}
\end{axis}
\end{tikzpicture}
\caption{Best IV RMSE per day of the calibrated sets embedded in the application (values above 150\,bp not shown). The gap in
2025 is a gap in the calibration campaigns, not in the market data.}\label{fig:history}
\end{figure}

%======================================================================
\section{Engineering}\label{sec:eng}
\begin{itemize}
  \item \emph{A static page}: \texttt{index.html} carries the page, its pricers and one day; the other days, the futures history
        and the calibrated sets are files loaded on demand. Settings are kept in the browser, or written to a folder of the disk
        with the File System Access API (Chrome, Edge). No server, no database, no account.
  \item \emph{Performance} (Chrome, Bates, six maturities, 222 options, COS): a full evaluation of the loss takes about
        $1.3$\,ms, a screen refresh about $3.5$\,ms. Characteristic-function constants that do not depend on $u$ are computed once
        per parameter set; complex arithmetic is written in real numbers in the hot loops; cosines are advanced by rotation
        instead of being recomputed.
  \item \emph{Mathematics on the page}: formulas are written in LaTeX and rendered once at build time with KaTeX, so the page
        runs no math script; rendered formulas are cached, and the daily GitHub Action needs only Python.
\end{itemize}

\section{Limits}
\begin{itemize}[nosep]
  \item The local optimisers depend on their starting point; differential evolution is global but far more expensive.
  \item Stacking days assumes parameters stable over the window; sharing $v_0$ is justified only for close days.
  \item Several compositions are weakly identified on a single day (\cref{sec:models}); a good RMSE does not mean that every
        parameter is meaningful.
  \item Writing files works only in browsers that implement the File System Access API.
\end{itemize}

\begin{thebibliography}{99}\small
\bibitem{alb} H.~Albrecher, P.~Mayer, W.~Schoutens and J.~Tistaert. The little Heston trap. \emph{Wilmott Magazine}, January 2007.
\bibitem{and} L.~Andersen. Simple and efficient simulation of the Heston stochastic volatility model. \emph{Journal of
  Computational Finance}, 11(3):1--42, 2008.
\bibitem{bn} O.~E.~Barndorff-Nielsen. Normal inverse Gaussian distributions and stochastic volatility modelling.
  \emph{Scandinavian Journal of Statistics}, 24(1):1--13, 1997.
\bibitem{bates} D.~Bates. Jumps and stochastic volatility: exchange rate processes implicit in Deutsche mark options.
  \emph{Review of Financial Studies}, 9(1):69--107, 1996.
\bibitem{cgmy} P.~Carr, H.~Geman, D.~Madan and M.~Yor. The fine structure of asset returns: an empirical investigation.
  \emph{Journal of Business}, 75(2):305--332, 2002.
\bibitem{chj} P.~Christoffersen, S.~Heston and K.~Jacobs. The shape and term structure of the index option smirk: why multifactor
  stochastic volatility models work so well. \emph{Management Science}, 55(12):1914--1932, 2009.
\bibitem{dps} D.~Duffie, J.~Pan and K.~Singleton. Transform analysis and asset pricing for affine jump-diffusions.
  \emph{Econometrica}, 68(6):1343--1376, 2000.
\bibitem{fo} F.~Fang and C.~Oosterlee. A novel pricing method for European options based on Fourier-cosine series expansions.
  \emph{SIAM Journal on Scientific Computing}, 31(2):826--848, 2008.
\bibitem{gh} F.~Gao and L.~Han. Implementing the Nelder--Mead simplex algorithm with adaptive parameters.
  \emph{Computational Optimization and Applications}, 51(1):259--277, 2012.
\bibitem{heston} S.~Heston. A closed-form solution for options with stochastic volatility with applications to bond and
  currency options. \emph{Review of Financial Studies}, 6(2):327--343, 1993.
\bibitem{kirkby} J.~L.~Kirkby. Efficient option pricing by frame duality with the fast Fourier transform. \emph{SIAM Journal on
  Financial Mathematics}, 6(1):713--747, 2015.
\bibitem{kou} S.~Kou. A jump-diffusion model for option pricing. \emph{Management Science}, 48(8):1086--1101, 2002.
\bibitem{ls} F.~Longstaff and E.~Schwartz. Valuing American options by simulation: a simple least-squares approach.
  \emph{Review of Financial Studies}, 14(1):113--147, 2001.
\bibitem{mcc} D.~Madan, P.~Carr and E.~Chang. The variance gamma process and option pricing. \emph{European Finance Review},
  2:79--105, 1998.
\bibitem{merton} R.~Merton. Option pricing when underlying stock returns are discontinuous. \emph{Journal of Financial Economics},
  3:125--144, 1976.
\bibitem{sp} R.~Storn and K.~Price. Differential evolution -- a simple and efficient heuristic for global optimization over
  continuous spaces. \emph{Journal of Global Optimization}, 11:341--359, 1997.
\end{thebibliography}
\end{document}
